\documentclass[10pt]{beamer}
\usepackage{csquotes}% Recommended

\usepackage{graphicx}
\usepackage{xcolor}
\usepackage{amsfonts}
\usepackage{amsmath}
\usepackage{amssymb}
\usepackage[skip=10pt plus1pt] {parskip}


%Information to be included in the title page:
\title{A Functional Overview of Integration}
\author{UChicago Political Science Math Camp}
\institute{Ryan Gibbons \& Luke Cain}
\date{2025}

\begin{document}

\frame{\titlepage}

\begin{frame}{Some Thoughts On Integration}
    \begin{itemize}
        \item Integration is the inverse of derivation, but can be conceptually more challenging
        \item Integral rules are often more difficult than derivative rules
        \item This session is more concerned with how to understand integrals than with how to compute them (i.e., intuition over execution)
        \item ``Spaghetti at the wall" pedagogical approach
    \end{itemize}
\end{frame}

\begin{frame}{Expectation}
    When calculating expected outcomes, we employ the following formula:

    \[E[u] = \sum_i \Pr[i] \cdot u_i\]

    The \textit{utility} of an event, then, is its value times its probability. \\

    \vspace{.5cm}

    Let $Heads = 1$ and $Tails = 0$. What is the expected value of a coin toss?
\end{frame}

\begin{frame}{More Expectation}
    We could model the expected roll of a D4 as

    \[\sum_i \Pr[i] \cdot u_i = \frac{1}{4} \cdot 1 + \frac{1}{4} \cdot 2 + \frac{1}{4} \cdot 3 + \frac{1}{4} \cdot 4 = 2.5\]

    What about a D6?\\
    \vspace{.5cm}
    \pause
    A D10? \\
    \vspace{.5cm}
    \pause
    A D100?
\end{frame}

\begin{frame}{Something Weird About Probability}
    What if we keep finding larger (fair) dice? What is $\lim_{i \to \infty}\Pr[x = i]$? \\

    \vspace{.5cm}
    \pause
    Assume that every outcome is equally likely (uniform distribution). We can evaluate:

    \[\sum_i^\infty \Pr[i] \cdot u_i = \sum_i^\infty \frac{1}{\infty} u_i\]

    And our summation and series rules lend that

    \[\implies \frac{1}{\infty} \cdot \sum_i^\infty u^i \implies \text{undefined}\]

    Something's not right here...
    
\end{frame}

\begin{frame}{A Qualifier}
    Let's take away the dice's fairness add a new condition: \\

    \vspace{.5cm}

    \textit{Every face of the dice has a unique value in $[-b,b]$} \\

    \vspace{.5cm}

    Then, as $n_{face} \to \infty$, every $x \in [-b,b]$ has a corresponding dice roll.

\end{frame}

\begin{frame}{Calculus and Infinity}

Calculus fundamentally solves the following contradictions: \\
\vspace{1cm}
\begin{itemize}

    \item Derivation solved for the rate of change of a function as $\Delta x \to 0$, such that $\frac{rise}{run} \to \infty$.
    
    \item Integration solves for the sum of an infinite number of infinitely small segments, where the size of each approaches $0$ such that $\sum_i \to 0$.

\end{itemize}

\textcolor{red}{In most cases, $\frac{rise}{run} \neq \infty$ and $\sum_i \neq 0$.}

\end{frame}

\begin{frame}{A Different Approach to Intuition}

    Say we want to measure the area bounded by some function(s), or the \textbf{area beneath the curve}:

    \begin{enumerate}
        \item If all functions are linear, we can use geometric formulas (extensions of $A = bh$) to evaluate
        \item If \underline{even one} function is non-linear\footnote{At least, non-circular, although the area of a circle is derived from calculus anyways}, then standard geometry can only approximate
    \end{enumerate}

\end{frame}

\begin{frame}{Integration From Approximation}
    We can approximate the area under a curve by splitting the area into a series of rectangles, such that $A = \sum_i bh_i$ \\

    \vspace{.5cm}

    With very large rectangles (i.e., very large $b$-values), this creates a lot of error. \textcolor{purple}{\href{https://www.desmos.com/calculator/fgqivwklhm}{But what happens}} as we add more rectangles, and shrink the size of each? \\

    \[A=\sum_i bh_i \quad \quad \implies \quad \quad b \to 0, i \to \infty\]

    \textcolor{red}{How does this reflect our derivation definition?}
\end{frame}

\begin{frame}{Notation}
    Integrate \textcolor{blue}{$f(x)$} with respect to \textcolor{red}{$x$} $\rightarrow$ What is the sum of every \textcolor{blue}{$f(x)$} over every possible \textcolor{red}{$x$}? \\

    \vspace{.5cm}

    We denote this as

    \[\int f(x) dx\]

    We call this an \textbf{indefinite integral}
\end{frame}

\begin{frame}{A New Limit Definition}
    \[\int f(x) dx = \lim_{b \to 0} \sum_i^{\frac{1}{b}}b f(x_i)\]

\end{frame}

\begin{frame}{Integration vs Derivation}
    \begin{itemize}
        \item The integral of $f(x) dx$ represents the sum of all $f$'s, such that if $f$ represents a rate of change, then $\int f(x) dx$ represents the \textbf{change that occurred...}
        \item Think: what function's derivative returns $f(x)$? That function represents $\int f(x)dx$
    \end{itemize}

To borrow again from Physics:

\[\text{\textbf{Jerk}} \Rightarrow \int{dt} = \text{\textbf{Acceleration}} \Rightarrow \int{dt} = \text{\textbf{Speed}} \Rightarrow \int dt = \text{\textbf{Position}}\]

\end{frame}

\begin{frame}{Definite vs. Indefinite Integrals}
    Often, we only want to evaluate the integral of a function over some specific domain. To integrate $f(x)$ w.r.t $x$ over the domain $[a,b]$, we write the \textbf{definite integral}

    \[\int_a^b f(x) dx\]

    \textcolor{red}{An indefinite integral often produces a function. A definite integral almost always returns a number (or function independent of $x$).}
\end{frame}

\begin{frame}{FTOC}

\[\int_a^b f'(x)dx = f(b) - f(a)\]

This yields two truths:

\begin{enumerate}
    \item The integral is the inverse of the derivative (the integral of a function's derivative is just the function)
    \item For some defined domain $[a,b]$, the value of the integral over that domain is $\int f(b)dx - \int f(a)dx$
\end{enumerate}

\end{frame}

\begin{frame}{FTOC and Indefinite Integrals}
    Let's reiterate an important point here: \\

    \vspace{.5cm}

    If we denote the indefinite integral $\int f(x)dx = F$, then $\int_a^b f(x)dx = F(b) - F(a)$ \\

    \vspace{.5cm}

    Then, to solve $\int_a^b f(x)dx$, we evaluate:

    \[\int f(x) dx = F(x) \quad \implies \quad \int_a^b f(x)dx = F(x) \bigg|_{x=a}^{x=b}\]
\end{frame}

\begin{frame}
    \thispagestyle{empty}
    \vspace{1cm}
    \begin{center}
        \begin{Huge}
            How Do We Actually Do This?
        \end{Huge}
    \end{center}
\end{frame}


\begin{frame}
    \begin{center}
        \includegraphics[width=5cm]{newton.jpeg}
    \end{center}
\end{frame}

\begin{frame}{The Good Stuff}
    We can use a series of rules to integrate a function in its general form \textit{\underline{for any value of $x$}}:

    \[f(x) \implies \text{[RULE]} \implies \int f(x) dx\]

    Note that now, our indefinite integral is itself another function of $x$.
\end{frame}

\begin{frame}{The Sum Rule}
    For all integrals:

    \[\int \Big[f(x) + g(x)\Big] dx  = \int f(x) dx + \int g(x) dx\]

    For definite integrals such that $a \leq b \leq c$:

    \[\int_a^c f(x) dx = \int_a^b f(x) dx + \int_b^c f(x)dx\]
\end{frame}

\begin{frame}{The Power Rule, Backwards}
    Recall the derivative power rule:

    \[\frac{d}{dx} x^k = k \cdot x^{k-1}\]

    Then, if $\int f'(x) dx = f(x)$, what function would, when derived, produce ours? \pause

    \[\int x^k dx = \frac{x^{k+1}}{k+1}\]

    plus one little thing...
\end{frame}

\begin{frame}{Remember Our Constants?}
    \begin{enumerate}
        \item We know that the derivative of a constant is $0$, such that $\frac{d}{dx} c = 0 \; \forall \; c \in \mathbb{R}$
        \item If integration ``reverts" a derivative, then it must account for any constant terms zero-d out in derivation
        \item We resolve this as follows:

        \[\int x^k dx = \frac{x^{k+1}}{k+1} + c\]

         \textit{even though we don't know what $c$ actually is.}
    \end{enumerate}
\end{frame}

\begin{frame}{A Power Rule Corollary (Reversed)}
    \begin{enumerate}
        \item We know that $\int x^k dx = \frac{1}{k+1} x^{k + 1}$ (\textbf{Power Rule}) and that $\int [f(x) + g(x)] dx = \int f(x) dx + \int g(x) dx$ (\textbf{Sum Rule})
        \item We also know that for a constant $\ell$, $\ell \cdot n = \sum_1^\ell n$, or $n$ added $\ell$ times. 
        \item Then, the integral is the sum of $x^k$ added up $\ell$ times:
        \[\int \bigg[ \ell \cdot x^k\bigg] dx = \sum_1^\ell \int x^k dx\]
        \item We know that $\sum_1^n x = nx$, so
        \[\sum_1^\ell \int x^k dx = \ell \cdot \int x^k dx \quad \implies \quad \int \bigg[ \ell \cdot x^k \bigg] dx = \ell \cdot \frac{1}{k+1} \cdot x^{k + 1}\]
    \end{enumerate}

    $\implies$  \textbf{Coefficients}: The coefficient of a term is kept as a coefficient of that term's integral (i.e., multiply $\frac{1}{k+1}$ by the coefficient as well)
\end{frame}

\begin{frame}{Quick Note on Constants}
    Recall that a constant $m$ can be represented as $m \cdot x^0$. Then,

    \[\int mdx = \int m \cdot x^0 dx\]
    \[\implies m \cdot \int x^0 dx\]
    \[\implies m \cdot \frac{x^{0+1}}{0+1} + c= m \cdot \frac{x^1}{1} + c\]
    
    \[\therefore \int m dx = mx + c\]
\end{frame}

\begin{frame}{Some Examples}
    Let's try some simple examples:

    \[\int x^2 + 2x^3 dx\] \pause

    \[\int 2x^3 + \frac{1}{2} x^4  dx \] \pause

    \[\int 3 + \sqrt[3]{x} dx\] \pause

    \[\int x^{-3} + x^{-2} + x^{-1} dx\] 
\end{frame}

\begin{frame}{Closing Some Subtle Foreshadowing}
    The power rule doesn't work for $f(x) = x^{-1} = \frac{1}{x}$, as it solves:

    \[\int x^{-1} dx = \frac{1}{-1+1} x^{-1 + 1} = - \frac{x^0}{0}\]

    Thankfully, our natural logarithm solves this for us:

    \[\int \frac{1}{x} dx = \ln(x) + c\]
\end{frame}

\begin{frame}{Integrating Exponentials}
    We have a specific rule for integrating exponentials and logarithms:

    \[\int a^x dx  = \frac{a^x}{\ln a} + c\]

    \[\int \log_b x dx = \frac{x}{\ln b} (\ln b - 1) + c\]
\end{frame}

\begin{frame}{Definite Integration}
    So far, we've expressed $F(x)$ for every $f(x)$, both as functions of $x$, as \textit{indefinite integrals}. When we add bounds of integration, we simply use FTOC:

    \[\int_a^b f(x)dx = F(b) - F(a)\]

    \[\int_2^4 x^3 dx \quad = \quad \int x^3 dx \bigg|_2^4 \]

    \[\implies \frac{x^4}{4} + c \;\bigg|_2^4\]

    But, we don't know $c$...
\end{frame}

\begin{frame}{$c$ you later}
    There's a handy trick with definite integrals:

    \[\int f(x) dx = F(x) + c\]
    \[\implies \int_a^b f(x) dx = F(b) + c - \Big( F(a) + c\Big)\]
    \[\implies \int_a^b f(x) dx = F(b) - F(a)\]

    For definite integrals, we subtract $c-c$ and can therefore exclude it. So,

    \[\frac{x^4}{4} + c \; \bigg|_2^4 \quad = \quad \frac{4^4}{4} - \frac{2^4}{4} \quad = 16\]
\end{frame}

\begin{frame}{More Examples}
    \[\int_0^2 x^3 + 2x^3 dx\] \pause

    \[\int_2^6 \frac{1}{2} x^2 + 4 dx\] \pause 

    \[\int_0^b \sqrt{x} + \sqrt{2} dx\] \pause

    \[\int_{a}^{2} \frac{1}{x^2} dx\] \pause
\end{frame}

\begin{frame}{A Turn For The Worst}
    Two of our fundamental derivative rules -- the \textbf{chain rule} and the \textbf{product rule} -- are a bit more complicated to invert:

    \[\frac{d}{dx}f(u) = f'(u) \cdot u'\]

    \[\implies F(u) = \int f(u) \cdot u du...\]

    There's not a clear solution here
\end{frame}

\begin{frame}{U-Substitution}
    The inverse of the Chain Rule is called \textbf{integration by substitution}, or \textbf{``u-sub"}, and works as follows:

    \begin{enumerate}
        \item Express the integral in the form
        \[\int f\Big(g(x)\Big) g'(x) dx\]
        \item Assign $g(x) = u$ such that the integral can be rewritten:
        \[ \int f(u) du\]
        \item Evaluate the integral:
        \[\int f(u) du = F(u)\]
        \item Substitute back $u = g(x)$:
        \[\int f\Big(g(x)\Big) g'(x) dx = F\Big(g(x)\Big)\]
    \end{enumerate}
\end{frame}

\begin{frame}{Integration By Parts}
    \textbf{Integration by parts} inverts the product rule. We will not cover it in math camp, but you should be aware of the concept and ready to learn its process.
\end{frame}

\begin{frame}{Area of a Circle}
    We know that $a = \pi r^2$, but the formal definition of $\pi$ is only $\pi = \frac{c}{d}$. \\

    \vspace{.5cm}

    Using \textit{only} this definition for $\pi$, how might we calculate the area of a circle?
\end{frame}

\begin{frame}{The Math}
    A circle can be divided into ``slices," and each slide can be approximated by a sum of triangles: $a_{slice} = \sum bh$. The error of these triangles approaches zero as we add more and more triangles... \\
    Imagine we take the sum of the area of an infinite number of triangles, rotated around the center of a circle. Then, since our triangle area is $a = \frac{1}{2}bh$, we have
    \[a_{circle} = \sum_i^\infty \frac{1}{i} \cdot \frac{1}{2}h\]
    The ``height" in this case is the radius of the circle. Now we have:
    \[\sum_i^\infty \frac{1}{i}\frac{1}{2} r = \int_{circle}\frac{1}{2}rdi = \frac{1}{2}r \cdot \int_{circle} di\]
    We know that $\int di = i$, so we just need to choose our bounds. We are integrating over the total distance of the circumference, which we know to be $2r \cdot \pi$. So, we can place the bounds:
    \[\int_0^{2r \cdot \pi}\frac{1}{2}rdi = \frac{1}{2}ri \; \bigg|_0^{2r \cdot \pi} = \frac{1}{2}r (2\pi r) - 0 = \pi r^2\]
\end{frame}

\begin{frame}{Improper Integrals}
    For definite integrals which do not converge to $f(x) = a$, or for any asymmetric integrals whereby $f(x) = \pm \infty$ for $x \in \mathbb{R}$, we call that integral \textbf{improper.} \\

    \vspace{.5cm}

    This can either result in a finite value, such as $\int_{-\infty}^{\infty} \frac{1}{1+x^2}dx$, or an undefined value, such as $\int_{-\infty}^{\infty} \frac{1}{x} dx$
\end{frame}

\begin{frame}{Gabriel's Horn}
    Imagine a horn formed by the function $f = \frac{1}{x}$ for the domain $x \in [1, \infty)$

    What is its volume? What is its surface area?
\end{frame}

\begin{frame}{Horn Volume}
    We know that $V = a \cdot h$; in this case we can set $h = \infty$ and $a = \pi r(x)^2 =\frac{\pi}{x^2}$. The volume is then the sum of this area over every possible height:

    \[V = \int_1^\infty \frac{\pi}{x^2} dx = - \frac{\pi}{x} \bigg|_1^{\infty} \]

    \[V = \pi\]

    The volume of the horn is finite.
\end{frame}

\begin{frame}{Horn Surface Area}
    We know that surface area can be calculated as $s = c \cdot h$. In this case, we set $h = \infty$ and $c = \pi 2r = \frac{2\pi}{x}$. The surface area is then the sum of this circumference over every possible diameter:

    \[s = \int_1^{\infty} \frac{2 \pi}{x} dx = 2\pi \cdot \ln(x) \bigg|_1^{\infty}\]

    \[s = \infty\]

The horn has finite volume but infinite surface area. \textit{You could fill the inside with paint, but never fully paint the outside.}
\end{frame}




\end{document}